\newpage

\section{Корреляционный анализ}

Для количественной оценки связи признаков с успехом стартапа был проведен корреляционный анализ (Рис.~\ref{fig:corr}). Бинарная переменная \texttt{success} (1 = success, 0 = failure) была сопоставлена с числовыми признаками, не подверженными ошибке выжившего: \texttt{is\_solo\_founder}, \texttt{tag\_count} и \texttt{desc\_length}. Для расчета использовался коэффициент корреляции Пирсона~\cite{pearson-corr}.

\begin{figure}[htbp]
    \centering
    \includegraphics[width=0.95\textwidth]{correlation_matrix.png}
    \caption{Матрица корреляций и корреляция признаков с успехом}
    \label{fig:corr}
\end{figure}

Результаты показали следующие значения корреляции с успехом:
\begin{itemize}
    \item \texttt{is\_solo\_founder}: $r = -0.11$ --- отрицательная корреляция, подтверждающая, что соло-основатели менее успешны
    \item \texttt{desc\_length}: $r = 0.09$ --- слабая положительная корреляция, согласующаяся с наблюдением о преимуществе умеренно длинных описаний
    \item \texttt{tag\_count}: $r = 0.04$ --- очень слабая положительная корреляция, объяснимая нелинейным характером зависимости
\end{itemize}

Между самими признаками наблюдается умеренная корреляция \texttt{tag\_count} и \texttt{desc\_length} ($r = 0.24$). Действительно, компании с более подробным описанием склонны указывать больше тегов. Признак \texttt{is\_solo\_founder} слабо отрицательно коррелирует с \texttt{tag\_count} ($r = -0.10$), что может свидетельствовать о том, что основатели-одиночки чаще фокусируются на узкой нише.

Важно отметить, что линейная корреляция Пирсона не улавливает нелинейные зависимости, выявленные в гипотезах H4 и H5. Перевернутая U-образная зависимость между количеством тегов и успехом проявляется как слабая линейная корреляция ($r = 0.04$), хотя реальная связь существенна. Это подчеркивает важность визуального анализа в дополнение к числовым метрикам.
